% Section 1 — Introduction
% Page budget: ~1 page (~720 words in acmart sigconf two-column).
% Every claim cross-checked against DOCS/docs8/PAPER-FINDINGS.md +
% DOCS/docs8/RQ-CLOSURE-TABLE.md + per-section SUMMARY.md files.

\section{Introduction}
\label{sec:intro}

A modern microservice cluster emits structured logs, distributed traces, numeric metrics, and Kubernetes events at the rate of tens of thousands of records per minute. Most of this stream is uneventful; a small fraction matters; a smaller fraction yet becomes a Jira ticket in the team's institutional memory of past incidents. For the engineer on-call, detecting that something is wrong is a solved problem---histogram baselines and learned classifiers over numeric features perform near-perfectly on \emph{``is this anomalous?''} The unsolved problem is what comes next: deciding which past Jira ticket describes the current failure. Today this lookup is performed by manual keyword search, takes ten to thirty minutes per page, and depends on the engineer recalling the right vocabulary.

Automating the lookup is harder than it appears. No public dataset jointly contains rich telemetry and linked Jira issues: the largest log corpora (HDFS, BGL, Thunderbird) have no ticket-side labels, the largest Jira archive (TAWOS, $458{,}232$ issues across $39$ open-source projects) has no telemetry, and production data from cooperating organizations is not shareable for legal and competitive reasons. Even when paired data exists, today's diagnostic systems sit at two design extremes. One runs every retriever and language-model verifier on every window---accurate but expensive in wall time, LLM tokens, and dollars. The other fixes a single pipeline and ignores the variety of evidence different incidents present: a CPU saturation looks different in metrics, a cascading Kafka outage looks different in traces, a third-party authentication blip looks different in logs. The choice between cost and coverage is made once at deployment, not per window.

This paper presents a capability-adaptive incident-triage agent that makes this choice once \emph{per window}. The agent inspects the available evidence for each window---which modalities are present, at what richness---and emits a per-window plan that invokes only the diagnostic methods the evidence warrants. The cheap path runs first; expensive retrievers and the language-model verifier escalate only when the cheap path is uncertain. An active evidence-gathering loop can fetch four additional signals from raw telemetry (Kubernetes events, expanded trace windows, Prometheus snapshots, peer-incident tickets) when retrieval consensus is low. Every diagnostic invocation is cacheable; every decision is replayable from a per-window execution trace that records the plan, the gates that fired, the skill outputs, and the budget spent.

We evaluate the agent on three datasets chosen to stress different parts of its design. The first is a controlled synthetic microservices benchmark of $1{,}008$ test windows, built from Google's Online Boutique demo with twenty-seven incident scenario families and a 347-ticket humanized Jira memory corpus. The second is the OpenTelemetry Demo, a structurally different polyglot application with $247$ test windows. The third is World of Logs, a corpus of $2{,}000$ real Apache Jira tickets across seven projects---including Kafka and MariaDB---with $304$ test windows. Every headline number carries a $1{,}000$-resample paired bootstrap confidence interval at seed $42$.

\paragraph{Contributions.}
\begin{itemize}
  \item \textbf{A capability-adaptive incident-triage agent} (\S\ref{sec:approach}). The agent emits a per-window diagnostic plan over eleven retrieval and composition primitives, four on-demand evidence-fetching tools, and a re-ranking step, gated by the evidence each window presents. Adding a new diagnostic method or evidence modality is a registry edit, not an architectural change.
  \item \textbf{Cost savings without retrieval-quality loss} (\S\ref{sec:results}). On synthetic microservices telemetry the agent matches the always-on cascade's Hit@5 of $0.758$ while spending $64.1\%$ less in dollar-equivalent LLM cost; the savings are robust to a tenfold sweep of the cheap-path confidence threshold.
  \item \textbf{Transfer to real production tickets} (\S\ref{sec:cross-app}). On real Apache Jira data the agent achieves Hit@1 $= 0.782$ and triage accuracy $= 0.934$, the highest values across all three datasets and, to our knowledge, the first agentic-retrieval result reported on a real Jira corpus paired with telemetry-style queries. The agent's Hit@5 dominates a BM25-only baseline by $1.4\times$ on real data and $8.7\times$ on synthetic data.
  \item \textbf{Statistically grounded negative findings} (\S\ref{sec:discussion}). The active evidence-gathering loop's positive Hit@1 contribution is not statistically distinguishable from zero (paired bootstrap $p = 0.18$); the extended-trace tool is significantly harmful on synthetic data ($p = 0.002$); and the agent's largest blind spot---a $40\%$ false-novelty rate on synthetic data---is induced by the language-model verifier and collapses by sixty-fold when that verifier is structurally disabled. These are findings the framework's instrumentation makes detectable; we report them rather than mask them.
  \item \textbf{A reproducible benchmark across three datasets}. Per-window traces, per-skill cost tables, paired-delta bootstrap confidence intervals, and a tool-use failure-mode catalog are emitted for every measurement. Replication is a single command per dataset.
\end{itemize}

The remainder of the paper is organized as follows. Section~\ref{sec:background} fixes terminology and surveys the relevant literature. Section~\ref{sec:approach} describes the agent. Section~\ref{sec:evaluation} defines the evaluation methodology---datasets, splits, baselines, and statistical envelope. Section~\ref{sec:results} reports within-dataset results; Section~\ref{sec:cross-app} reports cross-dataset transfer. Sections~\ref{sec:discussion}--\ref{sec:related} discuss findings, threats to validity, and related work, before concluding in Section~\ref{sec:conclusion}.
